\documentclass[letterpaper]{article}

\usepackage{latexsym}
\usepackage[empty]{fullpage}
\usepackage{titlesec}
\usepackage{marvosym}
\usepackage[usenames,dvipsnames]{color}
\usepackage{verbatim}
\usepackage{enumitem}
\usepackage[hidelinks]{hyperref}
\usepackage{fancyhdr}
\usepackage[english]{babel}
\usepackage{tabularx}
\usepackage{ragged2e}
\input{glyphtounicode}

\usepackage{xcolor}
\definecolor{accent}{HTML}{17685C}
\definecolor{accentdark}{HTML}{0E463E}
\hypersetup{
    colorlinks,
    linkcolor={accentdark},
    citecolor={accentdark},
    urlcolor={accentdark}
}
%----------FONT OPTIONS----------
\usepackage[default]{sourcesanspro}


\pagestyle{fancy}
\fancyhf{} % clear all header and footer fields
\fancyfoot{}
\renewcommand{\headrulewidth}{0pt}
\renewcommand{\footrulewidth}{0pt}

% Adjust margins
\addtolength{\oddsidemargin}{-0.5in}
\addtolength{\evensidemargin}{-0.5in}
\addtolength{\textwidth}{1in}
\addtolength{\topmargin}{-.6in}
\addtolength{\textheight}{1.2in}

\urlstyle{same}

\raggedbottom
\raggedright
\setlength{\tabcolsep}{0in}

% Sections formatting - small-caps heading with a thin underline rule
\let\oldsection\section
\renewcommand{\section}[1]{%
  \vspace{6pt}\noindent%
  {\large\scshape\textcolor{accentdark}{#1}}\par\vspace{-5pt}%
  {\color{accent}\rule{\textwidth}{1pt}}%
  \vspace{2pt}\par%
}

% Ensure that generate pdf is machine readable/ATS parsable
\pdfgentounicode=1

%-------------------------
% Vertical rhythm: standardize itemize spacing across all nesting levels
% so topics, sub-bullets, and consecutive items breathe consistently.
% Levels in this document:
%   itemize,1 = company/role list (label hidden)
%   itemize,2 = role's main bullets and topic items
%   itemize,3 = sub-bullets nested under a topic
\setlist[itemize]{leftmargin=1.4em, labelsep=0.4em,
                  parsep=0pt, partopsep=0pt}
\setlist[itemize,2]{itemsep=0.5pt, topsep=1pt}
\setlist[itemize,3]{itemsep=0.25pt, topsep=0.5pt, leftmargin=1.3em}

%-------------------------
% Custom commands
\newcommand{\resumeItem}[1]{%
  \item\small{\justifying #1}%
}

% Two-line role header rendered as a single tabular for perfectly uniform
% row spacing (no fragile negative \vspace between two tabulars).
\newcommand{\resumeSubheading}[4]{%
  \item%
    \begin{tabular*}{0.97\textwidth}{@{}l@{\extracolsep{\fill}}r@{}}%
      \textbf{\textcolor{accentdark}{#1}} & #2 \\%
      \textit{\small#3} & \textit{\small #4} \\%
    \end{tabular*}%
}

\newcommand{\resumeSubSubheading}[2]{%
    \item%
    \begin{tabular*}{0.97\textwidth}{@{}l@{\extracolsep{\fill}}r@{}}%
      \textit{\small#1} & \textit{\small #2} \\%
    \end{tabular*}%
}

\newcommand{\resumeProjectHeading}[2]{%
    \item%
    \begin{tabular*}{0.97\textwidth}{@{}l@{\extracolsep{\fill}}r@{}}%
      \small#1 & #2 \\%
    \end{tabular*}%
}

\newcommand{\resumeSubItem}[1]{\resumeItem{#1}}

% Topic sub-heading within a role: bold topic name, colon, then inline
% description sentence. Pair with a nested \resumeItemListStart/End if
% sub-bullets are needed under the topic.
\newcommand{\resumeTopic}[2]{%
  \item\small{\justifying \textbf{#1}: #2}%
}

% Sub-bullet marker (level 3): a small en-dash for a clean look.
\renewcommand{\labelitemiii}{\small\textendash}

\newcommand{\resumeSubHeadingListStart}{%
  \begin{itemize}[leftmargin=0.1in, label={}, itemsep=3pt, topsep=2pt]%
}
\newcommand{\resumeSubHeadingListEnd}{\end{itemize}}
\newcommand{\resumeItemListStart}{\begin{itemize}}
\newcommand{\resumeItemListEnd}{\end{itemize}}

%-------------------------------------------
%%%%%%  RESUME STARTS HERE  %%%%%%%%%%%%%%%%%%%%%%%%%%%%


\begin{document}

\begin{center}
    \textbf{\Huge \scshape \textcolor{accentdark}{Ashutosh Tiwari}} \\ \vspace{2pt}
    {\color{accent}\rule{\textwidth}{2pt}} \\ \vspace{2pt}
    \small 669-292-7534 $|$
    \href{mailto:findashutoshtiwari@gmail.com}{\underline{findashutoshtiwari@gmail.com}} $|$
    \href{https://linkedin.com/in/ashutosh--tiwari}{\underline{linkedin.com/in/ashutosh-{}-tiwari}} $|$
    \href{https://thunderock.github.io/}{\underline{thunderock.github.io}} $|$
    \href{https://github.com/thunderock}{\underline{github.com/thunderock}}
\end{center}


%-----------EXPERIENCE-----------
% \section{Summary}}
% Experienced Machine Learning Engineer with expertise in NLP, Graph ML, Large Language Models, and large-scale data pipelines. Proven ability to develop and optimize machine learning models, ensuring fairness and unbiased performance. Published researcher with skills in Python, PyTorch, and more. Committed to leveraging advanced ML techniques to drive innovation in fintech, health, and e-commerce sectors.

\section{Work Experience}
  \resumeSubHeadingListStart

  \resumeSubheading
    {Senior Machine Learning Engineer (ML Platform \& Frameworks)}{Jul. 2024 -- Present}
    {{\href{https://www.linkedin.com/company/adobe/}{\textcolor{accentdark}{\underline{ADOBE FIREFLY}}}}}{GENERATIVE AI, COMPUTER VISION, LARGE LANGUAGE MODELS}
    \resumeItemListStart
      \resumeItem{Lead fault-tolerant distributed inference and training infrastructure for the org's ML platform: GPU pipelines on Kubernetes engineered to survive node loss and partial failure, from hyperscaled daily EMR ingestion through production model inference.}
      \resumeTopic{Distributed Inference Frameworks}{Designed and led engineering for the org's online and offline inference systems powering experimentation and production at scale, reaching 1000s$\times$GPUs per model job across many model queues in one parallel wave.}
        \resumeItemListStart
          \resumeItem{Built online inference on vLLM + Ray for rapid experimentation with LLMs such as Falcon-40B and Llama~2~70B.}
          \resumeItem{Built offline inference on Ray Data + PyTorch Lightning; a module benchmarked on that engine at 1.09M rows, 8$\times$A100, 16 fractional-GPU actors. Introduced torch.compile(mode="max-autotune") in production: compiler-generated Triton kernels, fusion, CUDA-graph capture.}
          \resumeItem{Modeled pipelines as horizontally-scaling plugins with no pod-to-pod communication, KEDA-autoscaled on SQS queue depth; batched sends retry on exponential backoff with full jitter, and undeliverable work drains to a DLQ.}
        \resumeItemListEnd
      \resumeTopic{Foundation Model Training Framework}{Leading development of a PyTorch-based training framework adopted across the org for training generative models on billions of images and videos.}
        \resumeItemListStart
          \resumeItem{Designed a plugin-based strategy pattern over FSDP/FSDP2 to make distributed training strategies pluggable.}
          \resumeItem{Integrated Flash Attention 2/3, context parallelism, and distributed attention for DiT text-to-image and text-to-video training.}
          \resumeItem{Grew the engineering bench: mentored engineers onto the framework, authored its onboarding guides, and reviewed their integrations.}
        \resumeItemListEnd
      \resumeTopic{Build \& Deployment System}{Built a CLI-driven build system (BuildKit images on KEDA-scaled Kubernetes jobs, pushed to ECR), fronted by a Rust control plane and an MCP interface for AI-agent-driven builds, deployments, and inference.}
      \resumeTopic{Enrichment Service}{Orchestration over the inference framework; ran fire-and-forget pipelines behind a concurrency-safe S3 checkpoint cache, lease/TTL heartbeats and a circuit breaker.}
      \resumeTopic{Training-Data Governance}{Solely authored the org's fault-tolerant training-data lineage subsystem in Rust: idempotent claim/receipt registration, immutable S3 evidence layouts, and manifest-after-data publication.}
      \resumeTopic{Data Quality Framework}{Wrote the org's first data quality framework on Cerberus, used across enrichment pipelines to validate the correctness of feature-generation outputs across a Billion asset catalog growing 35--38M images/month.}
    \resumeItemListEnd

    \resumeSubheading
      {Software Dev Engineer II (ML Platform)}{Jan. 2019 -- Jul. 2021}
      {{\href{https://www.linkedin.com/company/swiggy-in/}{\textcolor{accentdark}{\underline{SWIGGY}}}}}{DATA SCIENCE PLATFORM, TIME SERIES FORECASTING}
      \resumeItemListStart
        \resumeTopic{Online Inference Platform}{Core contributor to the platform that served real-time model inference to 18M+ monthly transacting users across food delivery and quick commerce, autoscaled on live load metrics.}
          \resumeItemListStart
            \resumeItem{Engineered the dynamic request-routing layer that matched incoming API calls against geo-tags, traffic profile and service metadata, then dispatched them to model clusters partitioned by region and load domain.}
            \resumeItem{Implemented routing and dispatch on an Akka-style actor framework: resilient asynchronous invocation, fault isolation across clusters, and backpressure-aware scheduling that held sustained low P99 under high QPS.}
          \resumeItemListEnd
        \resumeTopic{Feature Store}{Built and maintained the feature store and its pipelines, serving on-demand features to production ML models at 4Bn rows and 10K QPS, with multichannel ingestion via Spark, Flink, and user files.}
        \resumeTopic{Forecasting \& Correlation Platform}{Founding member of the time-series forecasting platform adopted by teams across the org; its forecasts drove critical scaling decisions in real time.}
        \resumeTopic{DAQ}{Led DAQ, a large-scale API scraping tool collecting 15M rows daily for analysis and model training.}
      \resumeItemListEnd

    \resumeSubheading
      {Software Development Engineer (Search Relevance)}{Sep. 2017 -- Jan. 2019}
      {{\href{https://www.linkedin.com/company/flipkart/}{\textcolor{accentdark}{\underline{FLIPKART (a Walmart company)}}}}}{SEARCH RELEVANCE, QUERY INTENT, NLP}
      \resumeItemListStart
        \resumeTopic{Search Intent Models}{Built and improved CRF and neural-network intent models powering search and discovery for millions every day; drove the error-class analysis, designed the fixes and shipped them.}
        \resumeTopic{Tail Query Classifier}{Shipped a FastText-based query-store classifier that predicts the category of a tail query.}
        \resumeTopic{ML Workflow Automation}{Built the first workflow at Flipkart to retire manual deploys, automating training and auto-deployment of search models---written first in Luigi and later migrated to Airflow.}
          \resumeItemListStart
            \resumeItem{Wrote a generic Airflow framework that builds DAGs at runtime per ML model and orchestrates the train$\rightarrow$deploy flow, gating every auto-deploy on data and model validations.}
          \resumeItemListEnd
        \resumeTopic{Large-Scale Data Pipelines}{Implemented Cascading/HDFS pipelines over 4Bn+ user-event datapoints, extracting and transforming them into the training data the search models learned from.}
      \resumeItemListEnd

    \resumeSubheading
      {Software Development Engineer}{Sep. 2016 -- Sep. 2017}
      {{\href{https://www.linkedin.com/company/groupon/}{\textcolor{accentdark}{\underline{GROUPON}}}}}{BACKEND ENGINEERING, RECOMMENDATION SYSTEMS}
      \resumeItemListStart
        \resumeItem{\textit{Customer segmentation and deal-recommendation ML (NSVD unsupervised collaborative filtering on Neo4j); Cyclops customer-service platform live in every Groupon country.}}
      \resumeItemListEnd

    \resumeSubheading
      {Software Engineer}{Sep. 2015 -- Aug. 2016}
      {{\href{https://www.linkedin.com/company/netspeed-systems/about/}{\textcolor{accentdark}{\underline {NETSPEED SYSTEMS (Acquired by Intel)}}}}}{GRAPH ALGORITHMS, NETWORK ON CHIP}
      \resumeItemListStart
        \resumeItem{\textit{C++ graph algorithms for Network-on-Chip interconnects: Virtual Channel Arbitration, Polarity-based Arbitration, Multi-Cast Filtering, Structural Latency Breakdown.}}
      \resumeItemListEnd

  \resumeSubHeadingListEnd

\newpage
\section{Publications}

Accepted at \underline{\emph{\href{https://netsci2023.wixsite.com/netsci2023}{NetSci 2023}}} (Poster Presentation) and \underline{\emph{\href{https://www.ic2s2.org/}{IC2S2 2023}}} (Parallel Talk) as \textbf{first author}
\begingroup
\renewcommand{\section}[2]{}%
\begin{thebibliography}{}
 \bibitem{none} \textbf{Ashutosh Tiwari}, Prof. Sadamori Kojaku, Prof. Yong-Yeol Ahn,
\textquotedblleft Biased Contrastive Learning debiases Graph Neural Networks,\textquotedblright\ \textit{ International Conference on Network Science (NetSci)},
2023.\hfill \href{https://github.com/thunderock/residual2vec_/blob/master/paper/2023-07-03-fairness-NetSci2023.pdf}{\underline{Poster}}

A non-parametric contrastive framework for fairness-aware graph representation learning: it debiases GNN embeddings with respect to sensitive node attributes and structural homophily for more organic recommendations.
% \textbullet \quad Will be submitting full paper to \underline{\href{https://www.icwsm.org/2023/index.html/call_for_submissions.html}{ICWSM 2023}} as \textbf{first author} on \textbf{Contrastive Learning for Recommendation Systems using Graph Neural Networks} by May 15.
\end{thebibliography}
\endgroup
% \resumeItem{\textbullet \enskip Have submitted @\href{https://www.ic2s2.org/submit_abstract.html}{\underline{IC2S2}} as \textbf{first} author on \textbf{Fairness Aware Recommendation System and Graph ML}, pending decision notification}
% \resumeItemListEnd
%
%-----------EDUCATION-----------
\section{Education}
  \resumeSubHeadingListStart
    \resumeSubheading
      {Indiana University Bloomington}{}
      {Master of Science in Computational Data Science}{Aug. 2021 -- May 2023}
    \resumeSubheading
      {National Institute of Technology, Patna}{}
      {Bachelor of Science in Computer Science}{Aug. 2011 -- May 2015}
  \resumeSubHeadingListEnd
% https://www.linkedin.com/company/iuni/
%  Modelling research problems through networks
%
%Working on two projects:
%1. Working on embedding de-biasing novel training methods with Prof. YY Ahn and Prof. Sadamori Kojaku (as part of independent study for Fall '22)
%2. Working on user intent as a network and how that intent can be quantified. Working with Prof YY Ahn, Prof Preetesh Kantak and Prof Fahiz Baba Yara.  Project is funded by Kelly Businees School.
% \section{Publications}
% % \resumeItemListStart
% \resumeItem {Have submitted at \underline{\href{https://www2023.thewebconf.org/calls/research-tracks/fair-account-trans-ethic/}{WWW 23 Fairness Track}} as first author. Waiting for Acceptance Notification (Jan 2023)}.
% \resumeItemListEnd



\section{Research Experience / Independent Study}


    \resumeItemListStart
      \resumeItem {Debiasing independent study at \href{https://iuni.iu.edu/}{\underline{IUNI}} with Profs. YY Ahn \& S Kojaku --- novel training methods for fairness-aware graph recommendation from biased graphs (the line of work behind the NetSci/IC2S2 publication).}
      \resumeItem {Paid Research Assistant on \textbf{``User Intent as a Network''} with Profs. YY Ahn, P Kantak \& FB Yara, funded by the Kelly School of Business --- quantifying user intent through network structure.}
      \resumeItem {\href{https://nlp-lab.org/}{\underline{NLP Lab}}@IUB with Prof. D Cavar on \href{https://github.com/dcavar/tieml}{\underline{TieML}} and event-timeline modelling using fine-tuned large language models.}
     \resumeItemListEnd

\section{Teaching Experience}
    \resumeItemListStart
     \resumeItem {
     Teaching Assistant for \underline{\href{http://yongyeol.com/teaching/2023SP_netsci_syllabus.pdf}{Network Science}} (INFO-I 606) with Prof. YY Ahn}\hfill Spring 2023
     \resumeItem {
     Teaching Assistant for \underline{\href{https://cgi.luddy.indiana.edu/~rkhardon/B555-info.txt}{Machine Learning} }(CSCI-B 555) with Prof. R Khardon} \hfill Fall 2022

     \resumeItem {
     Teaching Assistant for \underline{\href{http://yongyeol.com/teaching/2022SP_netsci_syllabus.pdf}{Network Science}} (INFO-I 606) with Prof. YY Ahn}\hfill Spring 2022
     \resumeItemListEnd
%-----------PROJECTS-----------
\section{Selected Projects}
    \resumeSubHeadingListStart

\resumeProjectHeading
    {\textbf{rollout} $|$ \emph{Rust, Distributed Training \& Inference, RL for LLMs, vLLM, PyO3, Postgres}}{\href{https://github.com/thunderock/rollout}{\underline{github.com/thunderock/rollout}}}
    \resumeItemListStart
      \resumeItem{Multi-node reinforcement-learning framework for LLMs in Rust (18 crates, extensive test coverage, CI): PPO/GRPO/DPO/SFT/RM algorithms, vLLM-backed batch + online inference, actor/learner split with work-stealing, training-state and CRIU process snapshots, PyO3 in-process plugins.}
    \resumeItemListEnd
\resumeProjectHeading
    {\textbf{graph\_ml} $|$ \emph{Python, PyTorch Geometric, C++/Cython}}{\href{https://github.com/thunderock/graph_ml}{\underline{github.com/thunderock/graph\_ml}}}
    \resumeItemListStart
      \resumeItem{sklearn-style Graph-ML library on PyTorch/PyTorch Geometric with optimized C++/Cython random-walk kernels; CI-tested.}
    \resumeItemListEnd
\resumeProjectHeading
    {\textbf{bias\_manifold} $|$ \emph{PyTorch, Graph ML, Fairness}}{\href{https://github.com/thunderock/bias_manifold}{\underline{github.com/thunderock/bias\_manifold}}}
    \resumeItemListStart
      \resumeItem{Studies the structure and temporal evolution of bias manifolds across datasets --- the fairness-aware graph-representation line of work behind the NetSci/IC2S2 publication.}
    \resumeItemListEnd
\resumeProjectHeading
    {\textbf{BiasNet} $|$ \emph{Reinforcement Learning, PyTorch, Actor-Critic}}{\href{https://github.com/thunderock/BiasNet}{\underline{github.com/thunderock/BiasNet}}}
    \resumeItemListStart
      \resumeItem{Deep on-policy actor-critic RL agent that learns to play Street Fighter II from frame-difference observations (induced relational bias) over a large action space.}
    \resumeItemListEnd
\resumeProjectHeading
    {\textbf{DeepFoodie} $|$ \emph{Self-Supervised Learning, Embeddings}}{\href{https://github.com/thunderock/DeepFoodie}{\underline{github.com/thunderock/DeepFoodie}}}
    \resumeItemListStart
      \resumeItem{Self-supervised clustering of food items using learned ingredient embeddings.}
    \resumeItemListEnd
    \resumeSubHeadingListEnd

%
%-----------PROGRAMMING SKILLS-----------
\section{Technical Skills}
 \begin{itemize}[leftmargin=0.15in, label={}]
    \small{\item{
     \textbf{Languages \& Inference}{: Rust, C++, Python, PyO3, vLLM, ONNX Runtime, Triton, torch.compile, CUDA graphs, A100, H100} \\
     \textbf{Training \& Distributed}{: PyTorch, PyTorch Lightning, FSDP, Flash Attention, NCCL, Ray} \\
     \textbf{RL \& Alignment}{: PPO, GRPO, DPO, SFT, RLHF / reward modeling, actor-critic} \\
     \textbf{Platform \& Cloud}{: Kubernetes, KEDA, Docker / BuildKit, AWS (EMR, ECR, S3, SQS), Spark, Flink, Iceberg, Postgres} \\
     \textbf{Concepts}{: tensor parallelism, pipeline parallelism, Mixture of Experts (MoE) parallelism, KV-cache management, continuous batching} \\
     \textbf{ML \& Research}{: Graph Neural Networks, Contrastive Learning, Recommendation Systems, PyTorch Geometric, TensorFlow, scikit-learn, Pandas, NumPy}
    }}
 \end{itemize}
 % \section{Relevant Courses}}
 %  \begin{itemize}[leftmargin=0.15in, label={}]
 %     \small{\item{
 %      \textbf{Computational ML(B555), Deep Learning(E533), Reinforcement Learning(B659), Computer Vision(B657), Statistics(S520), Independent Study, Working on Fairness Aware AI and Bias manifolds in Graph ML (B669), Advanced Database Systems(B561)
 %      %\textbf{Computational Machine Learning(CSCI-B555), Deep Learning Systems(ENGR-E533), Applied ML(CSCI-P556), Computer Vision(CSCI-B657)
 % }{}
 %     }}
 %  \end{itemize}

%-------------------------------------------
\end{document}

